\section{The Witness of the Fathers, Doctors, and Liturgy}\label{sec:fathers-liturgy}

Near the end of the fourth century, preaching on the divinity of the Holy Spirit, St.\ Ambrose of Milan set down a single sentence that states both the height and the boundary of Catholic Marian theology: ``Mary was the temple of God, not the God of the temple. And therefore He alone is to be worshipped Who was working in His temple'' (\textit{On the Holy Spirit} III.11.80). To call the Mother of the Lord a temple, a shrine, or the living Ark of the Covenant never places a creature beside the Creator. It confesses that the eternal Word truly dwelt in her flesh, so that every honor paid to the living sanctuary passes through her and terminates in the God who made her his dwelling.

From the mid-second century to the nineteenth, Fathers, Ecumenical Councils, monastic preachers, scholastic Doctors, and the liturgical books of the Roman Church unfolded that conviction across Scripture. Rather than appearing overnight as a finished slogan, the Ark typology of Mary matured organically out of four deep roots: the second-century parallel between Eve and Mary, the fourth- and fifth-century confession of Mary as \textit{Theotokos} (``God-bearer''), the Christological and ecclesiological exposition of Psalm 131:8 and Revelation 11:19--12:1, and the praying liturgy of the Church.

\subsection{The New Eve, the God-Bearer, and the Guardrail of Worship (Second to Fourth Centuries)}

Long before preachers drew extended architectural parallels with Solomon's sanctuary, the earliest Fathers saw Mary's virginal obedience as the hinge on which salvation history turned. Writing around the middle of the second century, St.\ Justin Martyr pairs the virgin Eve's reception of the serpent's lie with the Virgin Mary's reception of Gabriel's word, anchoring the contrast directly in the overshadowing of Luke 1:35 and the consent of Luke 1:38 (\textit{Dialogue with Trypho} 100):

\begin{quote}
For Eve, who was a virgin and undefiled, having conceived the word of the serpent, brought forth disobedience and death. But the Virgin Mary received faith and joy, when the angel Gabriel announced the good tidings to her that the Spirit of the Lord would come upon her, and the power of the Highest would overshadow her: wherefore also the Holy Thing begotten of her is the Son of God; and she replied, ``Be it unto me according to thy word.''
\end{quote}

A generation later, St.\ Irenaeus of Lyons built that contrast into his theology of recapitulation. Because Adam and Eve had bound the human race to death through a virgin's unbelief, God willed that the knot be untied in reverse order through a virgin's faith. In \textit{Against Heresies} (III.22.4), Irenaeus traces that exact inversion:

\begin{quote}
In accordance with this design, Mary the Virgin is found obedient, saying, ``Behold the handmaid of the Lord; be it unto me according to thy word.'' But Eve was disobedient; for she did not obey when as yet she was a virgin. ... And thus also it was that the knot of Eve's disobedience was loosed by the obedience of Mary. For what the virgin Eve had bound fast through unbelief, this did the virgin Mary set free through faith.
\end{quote}

In the fifth book of the same work (\textit{Against Heresies} V.19.1), Irenaeus describes Mary's vocation with the Latin verb \textit{portare}---``to bear,'' ``to carry,'' ``to sustain,'' the very verb the Latin Bible uses of the Levites bearing the Ark of the Covenant:

\begin{quote}
For just as the former was led astray by the word of an angel, so that she fled from God when she had transgressed His word; so did the latter, by an angelic communication, receive the glad tidings that she should sustain (\textit{portaret}) God, being obedient to His word. And if the former did disobey God, yet the latter was persuaded to be obedient to God, in order that the Virgin Mary might become the patroness (\textit{advocata}) of the virgin Eve.
\end{quote}

For Irenaeus, Mary's bearing of God is never passive; she becomes the bearer of the Lord through free, believing obedience, so that Christ may take true human flesh from her and regenerate humanity from within.

In the fourth century, as the Church defended the full divinity and full humanity of the Son against Arianism and Apollinarianism, the title \textit{Theotokos} (``God-bearer,'' ``Mother of God'') moved to the center of orthodox preaching. St.\ Athanasius of Alexandria embeds the title inside his summary of the whole scriptural economy (\textit{Oration III against the Arians} 29):

\begin{quote}
Now the scope and character of Holy Scripture, as we have often said, is this,---it contains a double account of the Saviour; that He was ever God, and is the Son, being the Father's Word and Radiance and Wisdom; and that afterwards for us He took flesh of a Virgin, Mary Bearer of God, and was made man.
\end{quote}

In the same oration, Athanasius notes that Gabriel confessed himself sent by his Lord both to Zechariah and ``in the case of Mary, bearer of God'' (III.14), cites the Visitation directly---``John, while yet in the womb, leapt for joy at the voice of Mary Bearer of God''---and concludes that ``whereas the flesh is born of Mary Bearer of God, He Himself is said to have been born, who furnishes to others an origin of being'' (III.33).

A historical clarification belongs beside Athanasius's name. Devotional manuals frequently attribute to St.\ Athanasius a stirring Coptic encomium preserved on the Papyrus of Turin, saluting Mary as ``ark of the covenant, clothed with purity instead of gold,'' and similarly attribute Marian Annunciation homilies celebrating her as a golden ark to the third-century bishop St.\ Gregory Thaumaturgus. Modern patristic scholarship recognizes both texts as later pseudonymous compositions---a later Coptic pseudo-Athanasian homily and a fifth- or sixth-century pseudo-Gregorian corpus---rather than authentic third- or fourth-century works. Keeping those later homilies distinct from authentic fourth-century writings preserves the true historical growth of Ark typology: authentic fourth-century Fathers laid its Christological bedrock in \textit{Theotokos}, the Visitation, and the temple of the Spirit before later Byzantine and Coptic preachers elaborated explicit apostrophes to the Ark.

Across the Greek East, bishops repeated that same Christological bedrock. Instructing catechumens in Jerusalem around the middle of the fourth century, St.\ Cyril of Jerusalem gathers the testimonies to Christ's identity: ``The Father bears witness from heaven of His Son: the Holy Ghost bears witness, descending bodily in likeness of a dove: the Archangel Gabriel bears witness, bringing good tidings to Mary: the Virgin Mother of God bears witness: the blessed place of the manger bears witness'' (\textit{Catechesis} X.19, using Greek \textit{he Theotokos}). Writing to Cledonius against Apollinarius around 382---nearly fifty years before the Council of Ephesus---St.\ Gregory of Nazianzus states the rule of faith with sharp precision (\textit{Epistle 101}):

\begin{quote}
If anyone does not believe that Holy Mary is the Mother of God, he is severed from the Godhead. If anyone should assert that He passed through the Virgin as through a channel, and was not at once divinely and humanly formed in her (divinely, because without the intervention of a man; humanly, because in accordance with the laws of gestation), he is in like manner godless.
\end{quote}

Gregory's rejection of the ``channel'' image (\textit{dia solenos}) illuminates the supreme superiority of the living Ark over the Mosaic chest. The acacia box of Sinai merely contained stone tablets without contributing anything to their substance. Mary is not a hollow conduit through whom a heavenly body merely passed; the eternal Son was truly formed in her according to the laws of gestation, taking his very human flesh from hers.

Gregory's fellow Cappadocian, St.\ Gregory of Nyssa, brings the Sinai tablets---the very covenant deposit placed inside the Ark---directly into the mystery of the Annunciation. Commenting on the broken tablets of the Law in his \textit{Life of Moses} (transmitted in St.\ Thomas Aquinas's \textit{Catena Aurea on Luke} 1:35), Gregory writes:

\begin{quote}
For the tables of our nature which guilt had broken, the true Lawgiver has formed anew Himself from our dust without cohabitation, creating a body capable of taking His divinity, which the finger of God hath carved, that is to say, the Spirit coming upon the Virgin.
\end{quote}

In the Latin West, St.\ Ambrose of Milan united this Christological rigor with luminous sanctuary language. Commenting in \textit{On the Holy Spirit} (III.11.79--80) on Psalm 98:5 (``Exalt ye the Lord our God, and adore his footstool, for it is holy'') and Luke 1:35, Ambrose explains that the ``footstool'' adored in the psalm signifies the earth of Christ's flesh, ``which we this day also adore in the mysteries, and which the apostles, as we said above, adored in the Lord Jesus.'' Because that adorable flesh was formed by the Holy Spirit's overshadowing, the Spirit himself is rightly adored as God. Immediately Ambrose sets the permanent boundary separating the adoration owed to God alone (\textit{latria}) from the singular veneration paid to his Mother (\textit{hyperdulia}):

\begin{quote}
Since, then, the mystery of the Incarnation is to be adored, and the Incarnation is the work of the Spirit, as it is written, ``The Holy Spirit shall come upon thee, and the power of the Most High shall overshadow thee, and that Holy Thing Which shall be born of thee shall be called the Son of God,'' without doubt the Holy Spirit also is to be adored, since He Who according to the flesh was born of the Holy Spirit is adored. And let no one divert this to the Virgin Mary; Mary was the temple of God, not the God of the temple. And therefore He alone is to be worshipped Who was working in His temple.
\end{quote}

With adoration guarded strictly for the triune God, Ambrose speaks of Mary's consecrated person with unmistakable temple imagery. In his \textit{Exposition of the Gospel of Luke}, he writes that by calling Joseph a just man, Scripture ``sufficiently declared that he could not violate the temple of the Holy Spirit, the womb of the mystery, the mother of the Lord'' (\textit{sancti Spiritus templum, uterum mysterii, matrem Domini violare non potuit}, II.6). In the next section he draws out the bond between Mary and the Church: ``Well is she betrothed, yet a virgin; because she is a type of the Church, which is immaculate, yet wedded. A virgin conceived us by the Spirit; a virgin brings us forth without groaning'' (\textit{quia est Ecclesiae typus, quae est immaculata, sed nupta}, II.7). Contemplating the Visitation, he calls her ``pregnant with the Word, full of God'' (\textit{mater Domini Verbo foeta, Deo plena est}, II.25), and in \textit{Concerning Virgins} (II.2.15, 18) he praises her as ``the Virgin within the house, the companion for the ministry, the mother at the temple,'' whose soul, ``once more uncovered by the hand of the Eternal Priest, exhales the vapour of the divine fire.''

\subsection{Augustine, Ephesus, and the Maturing of Psalm 131:8 (Fifth to Eighth Centuries)}

At the opening of the fifth century, St.\ Augustine of Hippo reveals the exact theological soil out of which later Marian Ark exegesis grew. His exposition of Psalm 131:8 (Hebrew 132:8) is decisive. In the Septuagint (\textit{anastethi, Kyrie, eis ten anapausin sou, sy kai he kibotos tou hagiasmatos sou}), the Gallican Vulgate (\textit{Surge, Domine, in requiem tuam, tu et arca sanctificationis tuae}), and the Douay-Rheims Bible, that royal verse reads: ``Arise, O Lord, into thy resting place: thou and the ark, which thou hast sanctified''---diverging from the Masoretic Hebrew text (and from all versions of 2 Chronicles 6:41), which reads ``the ark of thy strength.''

Expounding Psalm 131:8 in his \textit{Expositions on the Psalms}, Augustine does not apply the verse directly to Mary's bodily Assumption. He unfolds first an ecclesiological sense and notes beside it an older Christological interpretation:

\begin{quote}
``Arise, O Lord, into Thy resting place.'' He saith unto the Lord sleeping, ``Arise.'' Ye know already who slept, and who rose again. ... ``Thou, and the ark of Thy sanctification:'' that is, Arise, that the ark of Thy sanctification, which Thou hast sanctified, may arise also. He is our Head; His ark is His Church: He arose first, the Church will arise also. The body would not dare to promise itself resurrection, save the Head arose first. The Body of Christ, that was born of Mary, hath been understood by some to be the ark of sanctification; so that the words mean, Arise with Thy Body, that they who believe not may handle.
\end{quote}

Augustine's commentary shows the organic trajectory of Catholic typology. First, Christian exegetes saw Christ's own human body, born of Mary, as the Ark of sanctification rising from death into eternal rest. Next, Augustine saw the whole Church, united to her risen Head as his Mystical Body, as the Ark rising after him. Because Mary is both the fleshly source of Christ's humanity and the singular archetype of the Church, later Eastern and Western preaching naturally beheld in her bodily Assumption the firstfruits of the Church's ascent into the Lord's resting place.

Elsewhere Augustine delights in the same paradox of the infinite God borne by a human mother. Preaching on the Nativity (\textit{Sermon} 185.1), he sets Psalm 84:12 alongside the manger: ``Truth, which is in the bosom of the Father, sprang from the earth so that He might be also in the bosom of a mother; Truth, by which the world is contained, sprang from the earth so that He might be borne in a woman's hands'' (\textit{Veritas qua mundus continetur, de terra orta est, ut femineis manibus portaretur}). And in \textit{On Nature and Grace} (36 [42]), answering Pelagius on personal sin, he sets Mary apart solely from her union with her Son: ``We must except the holy Virgin Mary, concerning whom I wish to raise no question when it touches the subject of sins, out of honour to the Lord; for from Him we know what abundance of grace for overcoming sin in every particular was conferred upon her who had the merit to conceive and bear Him who undoubtedly had no sin.''

While Augustine preached in Hippo, the Nestorian crisis in the East brought the title \textit{Theotokos} to solemn conciliar definition at the Council of Ephesus (431) and the Council of Chalcedon (451). In his \textit{Commentary on the Gospel of Luke} (Sermon I), St.\ Cyril of Alexandria states the exact union that makes Mary both Mother of God and living temple:

\begin{quote}
For a union of two natures was made, and therefore we confess One Christ, One Son, One Lord. And it is with reference to this notion of a union without confusion that we proclaim the holy Virgin to be the mother of God, because God the Word was made flesh and became man, and by the act of conception united to Himself the temple that He received from her.
\end{quote}

A century later, the Second Council of Constantinople (Fifth Ecumenical Council, 553) solemnly joined ``Godbearer'' (\textit{Theotokos}) and ``ever-virgin'' (\textit{aeiparthenos}), reaffirming Ephesus and Chalcedon against any attempt to reduce Mary to the mother of a mere man later joined to God. Canon VI declares:

\begin{quote}
If anyone says that the holy, glorious, ever-virgin Mary is called Godbearer by abuse and not truly, or by analogy, as though a mere man were born of her, and not as though God the Word were incarnate of her ... or, if anyone calls her manbearer or Christbearer, as though Christ were not God, and does not confess her as Godbearer, in the proper sense and in truth, because God the Word, who was begotten of the Father before all worlds, was incarnate of her in the last days; and (does not confess) that in this pious sense the holy Synod of Chalcedon confessed her to be Godbearer,---let him be anathema.
\end{quote}

Across the Christian East from the fourth through the eighth centuries---beginning with Syriac hymnography in St.\ Ephrem the Syrian and culminating in Byzantine Dormition preaching by St.\ Andrew of Crete, St.\ Germanus of Constantinople, and St.\ John of Damascus---the imagery of Exodus, the Tabernacle, and the Ark of sanctification flowered into liturgical poetry and homilies. St.\ John of Damascus, the great summarizer of Eastern patristic doctrine, anchors that sanctuary theology firmly in the Incarnation in his \textit{Exposition of the Orthodox Faith} (\textit{De fide orthodoxa} III.12):

\begin{quote}
God the Word Himself, Who was begotten of the Father timelessly before the ages, and was with the Father and the Spirit without beginning and through eternity, took up His abode in these last days for the sake of our salvation in the Virgin's womb, and was without change made flesh and born of her. For the holy Virgin did not bare mere man but true God: and not mere God but God incarnate, Who did not bring down His body from Heaven, nor simply passed through the Virgin as channel, but received from her flesh of like essence to our own and subsisting in Himself.
\end{quote}

Turning to Mary's preparation in the fourth book of the same treatise (\textit{De fide orthodoxa} IV.14), John contemplates her consecration as God's dwelling place:

\begin{quote}
Then planted in the House of God and increased by the Spirit, like a fruitful olive tree, she became the home of every virtue, turning her mind away from every secular and carnal desire, and thus keeping her soul as well as her body virginal, as was meet for her who was to receive God into her bosom: for as He is holy, He finds rest among the holy. Thus, therefore, she strove after holiness, and was declared a holy and wonderful temple fit for the most high God.
\end{quote}

In the same chapter he writes that she ``ministered to the Creator in that He was created, to the Fashioner in that He was fashioned, and to the Son of God and God in that He was made flesh and became man from her pure and immaculate flesh and blood, satisfying the debt of the first mother.'' In his three celebrated \textit{Homilies on the Dormition}, John of Damascus carried this living-Ark typology and the verse of Psalm 131:8 to classic Byzantine maturity on the feast of Mary's bodily passage into heaven, a patristic voice later invoked by Pope Pius XII in the Apostolic Constitution \textit{Munificentissimus Deus} (1950).

At the very same period on the northern edge of the Latin West, the Venerable Bede commented on the vision of the heavenly Ark and the Woman clothed with the sun in his \textit{Explanation of the Apocalypse} (Book II, on Rev 11:19--12:1). Centuries before medieval Bibles inserted a chapter break between those two verses, Bede read them as one continuous revelation of the Incarnation and the Church:

\begin{quote}
The temple of God, which was placed aforetime on the earth, had the ark of the covenant, inclosed within the mystical veil. But now, in the Church, which is ``the temple of the living God,'' whose ``conversation is in heaven,'' the ark of His Incarnation is laid open to the whole world, as the veil of the old temple, and the middle wall of partition, have been rent asunder by the blood of the Lord. For like the heavenly manna in clean gold is the Divinity in a holy body. ... \textit{sign.} It is the same sign which now also appears in the Church, that God is born of man. \textit{a woman.} The Church is clothed with the light of Christ, and treads upon temporal glory.
\end{quote}

Bede's simile---``like the heavenly manna in clean gold is the Divinity in a holy body''---recalls Hebrews 9:4's golden urn of manna inside the Ark and holds Christ, his birth from Mary (``that God is born of man''), and the Church together in a single vision.

\subsection{Monastic, Scholastic, and Early Modern Reception (Twelfth to Seventeenth Centuries)}

In the twelfth century, Cistercian monastic preaching clothed patristic sanctuary theology in poetic warmth. Preaching on the Annunciation (\textit{Super Missus est} II.7), St.\ Bernard of Clairvaux contemplates Gideon's fleece drenched with heavenly dew while the surrounding floor remains dry as a figure of the virginal conception:

\begin{quote}
This similitude represents the flesh assumed from the flesh of Mary without injury to her virginity. Upon it Heaven dropped down dew, filling it with the plenitude of the Divinity, and from that fullness we have all received---we who, but for it, were as parched and arid soil.
\end{quote}

Contemplating Christ's birth in his \textit{Sermon on the Vigil of the Nativity}, Bernard speaks of Mary's womb in the language of temple consecration: ``O birth alone without sorrow, alone without shame, free from corruption, not unlocking, but consecrating the temple of the Virgin's womb!''

A century later, in the \textit{Tertia Pars} of the \textit{Summa theologiae} (qq.\ 27--30, 35), St.\ Thomas Aquinas gathered a millennium of patristic temple and shrine language into a unified theology grounded in the hypostatic union. Because Christ is truly God incarnate, ``it can be truly said that God was conceived and born of the Virgin... Therefore it is heretical to deny that the Blessed Virgin is the Mother of God'' (III, q.\ 35, a.\ 4). From this maternal proximity to the Son flows Mary's unmatched fullness of grace (III, q.\ 27, a.\ 5):

\begin{quote}
In every genus, the nearer a thing is to the principle, the greater the part which it has in the effect of that principle... Now Christ is the principle of grace, authoritatively as to His Godhead, instrumentally as to His humanity: whence (John 1:17) it is written: ``Grace and truth came by Jesus Christ.'' But the Blessed Virgin Mary was nearest to Christ in His humanity: because He received His human nature from her. Therefore it was due to her to receive a greater fulness of grace than others.
\end{quote}

Aquinas then connects the dwelling of Divine Wisdom in Mary's womb directly to her complete immunity from actual sin (III, q.\ 27, a.\ 4):

\begin{quote}
God so prepares and endows those, whom He chooses for some particular office, that they may be rendered capable of fulfilling it... Thirdly, because of the singular manner in which the Son of God, who is the ``Divine Wisdom'' (1 Cor.\ 1:24) dwelt in her, not only in her soul but in her womb. And it is written (Wis.\ 1:4): ``Wisdom will not enter into a malicious soul, nor dwell in a body subject to sins.'' We must therefore confess simply that the Blessed Virgin committed no actual sin, neither mortal nor venial; so that what is written (Cant 4:7) is fulfilled: ``Thou art all fair, O my love, and there is not a spot in thee.''
\end{quote}

Defending Mary's perpetual virginity after Christ's birth against Helvidius (III, q.\ 28, a.\ 3), Aquinas invokes the ancient liturgical title \textit{Sacrarium Spiritus Sancti}, ``shrine of the Holy Ghost'':

\begin{quote}
Secondly, this error is an insult to the Holy Ghost, whose ``shrine'' was the virginal womb, wherein He had formed the flesh of Christ: wherefore it was unbecoming that it should be desecrated by intercourse with man. ... Fourthly, it would be tantamount to an imputation of extreme presumption in Joseph, to assume that he attempted to violate her whom by the angel's revelation he knew to have conceived by the Holy Ghost. We must therefore simply assert that the Mother of God, as she was a virgin in conceiving Him and a virgin in giving Him birth, did she remain a virgin ever afterwards.
\end{quote}

Finally, asking why God announced the Incarnation through Gabriel before conceiving his Son in Mary's womb (III, q.\ 30, a.\ 1), Aquinas quotes St.\ Augustine's \textit{On Holy Virginity}---``Mary is more blessed in receiving the faith of Christ, than in conceiving the flesh of Christ''---and adds the covenantal reason that unites Mary to the whole human family: ``Fourthly, in order to show that there is a certain spiritual wedlock between the Son of God and human nature. Wherefore in the Annunciation the Virgin's consent was besought in lieu of that of the entire human nature.''

Aquinas also illustrates how theological understanding can mature within fidelity to Christ. Accepting Mary's bodily Assumption as established doctrine (III, q.\ 27, a.\ 1), he nevertheless hesitated in III, q.\ 27, a.\ 2 to affirm her sanctification before animation, fearing that preservation from original sin might seem to exempt her from Christ's universal redemption. A generation later, Blessed John Duns Scotus resolved that difficulty in his commentary on Book III of the \textit{Sentences} (dist.\ 3) by distinguishing the order of nature from the order of time and showing that preservative redemption---preventing a fall before it occurs rather than raising one after it---is the most perfect act of a Redeemer, so that Mary owes her spotless beginning wholly to the foreseen merits of her Son.

In the sixteenth and seventeenth centuries, Catholic biblical scholars brought close attention to the Hebrew, Greek, and Latin texts together with patristic exegesis. In his \textit{Commentary on the Book of Psalms}, St.\ Robert Bellarmine expounds Psalm 131:7--8 with lucid architectural and textual accuracy:

\begin{quote}
The sanctuary in the interior of the tabernacle was so constructed that the Ark was the footstool, and the propitiatory was the seat of the Lord, and was supported by two Cherubim. ... And he lets us see that, when he speaks thus to the Lord, he does not address him, as he essentially exists, for, as such, ``the heaven of heavens do not contain him,'' but as he is peculiarly in the Ark of the Testament, because it was from it he gave his answers, and he, therefore, adds, ``Thou and the ark which thou hast sanctified;'' that is, you and your throne, and the footstool of your feet, which footstool is ``the ark which thou hast sanctified,'' and through which you will be sanctified and honored by all who know it. In the Hebrew it is the ``Ark of thy strength,'' because it was through it God displayed his strength.
\end{quote}

Bellarmine immediately recalls the dreadful sanctity of that Ark at Beth-shemesh and in the sudden death of Uzzah (Oza). In the same era, Cornelius a Lapide compiled patristic and medieval commentary across the Gospels, explaining the overshadowing veil of Luke 1:35, upholding \textit{Theotokos} from Elizabeth's cry in Luke 1:43, and citing St.\ Bernard on the humility of the \textit{Magnificat} in Luke 1:46. Yet when one reads both St.\ Ambrose (\textit{Exposition of the Gospel of Luke} II.29) and Cornelius a Lapide on Luke 1:56 (``about three months''), a revealing historical fact appears: both dwell richly on the sanctification of John and Elizabeth during Mary's three-month sojourn, while neither draws the detailed narrative comparison with David's procession and Obed-edom's house in 2 Samuel 6. The Church's conviction that Mary is the living Ark of the Covenant is patristic, medieval, and liturgical; the systematic narrative comparison of Luke 1 with 2 Samuel 6 is a modern development of canonical biblical theology deepening an inheritance already centuries old.

\subsection{The Church Praying: Liturgical Propers, the Litany of Loreto, and Ineffabilis Deus}

The Church's \textit{lex orandi}---her public worship and received prayer---places Mary directly inside the sanctuary texts of Exodus, Chronicles, the Psalms, and the Apocalypse. Both the 1962 \textit{Missale Romanum} and the postconciliar \textit{Missale Romanum} with its \textit{Ordo lectionum Missae} (1981) weave Ark and dwelling-place Scripture into the solemnities of the liturgical year, each according to its own calendar and selection of readings.

\subsubsection{The Solemnity of the Assumption (14--15 August)}

Nowhere does Ark typology shine more openly in the Roman rite than on the Vigil and Day of the Assumption.

In the 1962 \textit{Missale Romanum}, the Vigil (14 August) retains an ancient Collect addressing God directly in sanctuary language: \textit{Deus, qui virginalem aulam beatae Mariae, in qua habitares, eligere dignatus es}---``O God, who didst vouchsafe to choose the virginal hall of blessed Mary wherein to dwell.'' Its Offertory acclaims her who ``didst bear the Creator of all'' (\textit{Beata es, Virgo Maria, quae omnium portasti Creatorem: genuisti qui te fecit, et in aeternum permanes Virgo}), and its Communion sings, ``Blessed is the womb of the Virgin Mary, which bore the Son of the eternal Father'' (\textit{Beata viscera Mariae Virginis, quae portaverunt aeterni Patris Filium}), following Sirach 24:23--31 and Luke 11:27--28. For the Mass during the Day (15 August), promulgated in 1950 at the solemn definition of the Assumption, the 1962 Missal opens with Revelation 12:1 as its Introit (\textit{Signum magnum apparuit in caelo}, ``A great sign appeared in heaven: a woman clothed with the sun...''), reads Judith 13:22--25 and 15:10 as its Lesson, sings Genesis 3:15 at the Offertory, and proclaims the Visitation and \textit{Magnificat} (Luke 1:41--50) as its Gospel, while the Collect confesses that God assumed ``the immaculate Virgin Mary, Mother of thy Son, in body and soul to heavenly glory.''

In the postconciliar \textit{Ordo lectionum Missae} (1981), Ark texts form the opening readings of both Masses. The Vigil (14 August) takes as its First Reading 1 Chronicles 15:3--4, 15--16 and 16:1--2---King David assembling Israel and the Levites to carry the Ark of the Covenant on their shoulders with staves into the tent pitched in Jerusalem---paired with Psalm 132 (Vulgate 131):6--7, 9--10, 13--14, followed by 1 Corinthians 15:54b--57 and Luke 11:27--28. The Mass during the Day (15 August) retains Revelation 12:1 (\textit{Signum magnum}) as its Entrance Antiphon and opens its First Reading across the medieval chapter break at Revelation 11:19a---``And the temple of God was opened in heaven: and the ark of his testament was seen in his temple''---continuing directly into Revelation 12:1--6a and 10ab, followed by Psalm 45 (Vulgate 44), 1 Corinthians 15:20--27, and the complete Visitation and \textit{Magnificat} pericope, Luke 1:39--56.

\subsubsection{The Immaculate Conception, Annunciation, and Visitation}

The same sanctuary theology runs across the other principal Marian feasts of the Roman calendar.

On the Immaculate Conception (8 December), both the 1962 \textit{Missale Romanum} and the postconciliar \textit{Missale Romanum} share the classic Collect: \textit{Deus, qui per immaculatam Virginis Conceptionem dignum Filio tuo habitaculum praeparasti: quaesumus; ut, qui ex morte eiusdem Filii tui praevisa eam ab omni labe praeservasti, nos quoque mundos eius intercessione ad te pervenire concedas}---``O God, who by the Immaculate Conception of the Virgin didst prepare a worthy dwelling place for thy Son: we beseech thee that, as in view of the foreseen death of the same thy Son thou didst preserve her from all stain, so thou wouldst grant us also, cleansed by her intercession, to come unto thee.'' The 1962 Missal pairs Proverbs 8:22--35 with Luke 1:26--28; the postconciliar lectionary pairs Genesis 3:9--15, 20 and Ephesians 1:3--6, 11--12---bringing St.\ Paul's ``in which he hath graced us'' (\textit{echaritosen}, Eph 1:6) directly beside Gabriel's ``full of grace'' (\textit{kecharitomene}, Luke 1:28)---with the entire Annunciation Gospel, Luke 1:26--38.

On the Annunciation (25 March), both books pair Isaiah 7 (the virgin conceiving Emmanuel) with Luke 1:26--38. The 1962 Collect prays that we ``who truly believe her to be the Mother of God'' (\textit{qui vere eam Genetricem Dei credimus}) may be assisted by her intercessions; the postconciliar solemnity, titled \textit{The Annunciation of the Lord}, adds Hebrews 10:4--10 (``Sacrifice and oblation thou wouldest not: but a body thou hast fitted to me,'' Heb 10:5) as its Second Reading.

On the Visitation (kept on 2 July in the 1962 Missal and on 31 May in the postconciliar calendar), the 1962 Mass pairs Song of Songs 2:8--14 (``behold he cometh leaping upon the mountains, skipping over the hills'') with Luke 1:39--47, and its Secret confesses that the humanity of the Only-Begotten, born of the Virgin, ``did not diminish his mother's integrity, but consecrated it'' (\textit{matris integritatem non minuit, sed sacravit}). The postconciliar Feast sets Zephaniah 3:14--18 (``Give praise, O daughter of Sion... the king of Israel, the Lord, is in the midst of thee'') or Romans 12:9--16b before the complete Visitation narrative, Luke 1:39--56.

\begin{fourstable}{Feast}{1962 Missale Romanum}{Postconciliar Missal / 1981 Lectionary}{Sanctuary and Ark Texts}
Assumption Vigil (14 Aug) & Sir 24:23--31; Luke 11:27--28; Collect: \textit{virginalem aulam... in qua habitares} & 1 Chr 15:3--4, 15--16; 16:1--2; Ps 132(131):6--7, 9--10, 13--14; 1 Cor 15:54b--57; Luke 11:27--28 & Virginal hall of God's dwelling (1962); Levitical transfer of the Ark and Zion Ark psalm (postconciliar)\\
Assumption Day (15 Aug) & Introit Rev 12:1; Jdt 13:22--25, 15:10; Offertory Gen 3:15; Luke 1:41--50 & Entrance Rev 12:1; Rev 11:19a, 12:1--6a, 10ab; Ps 45(44); 1 Cor 15:20--27; Luke 1:39--56 & Woman clothed with the sun and Visitation (both); heavenly Ark of Rev 11:19a joined directly to Rev 12:1 (postconciliar)\\
Immaculate Conception (8 Dec) & Prov 8:22--35; Luke 1:26--28; Collect: \textit{dignum Filio tuo habitaculum praeparasti} & Gen 3:9--15, 20; Eph 1:3--6, 11--12; Luke 1:26--38; shared Collect (\textit{habitaculum}) & Worthy dwelling place prepared by preservative grace (shared); Eph 1:6 joined to Luke 1:28 (postconciliar)\\
Annunciation (25 Mar) & Isa 7:10--15; Ps 44(45); Luke 1:26--38; Collect confessing \textit{Genetricem Dei} & Isa 7:10--14, 8:10; Ps 40(39); Heb 10:4--10; Luke 1:26--38 & Virgin conceiving Emmanuel and overshadowing of Luke 1:35 (both); body prepared for Christ (postconciliar)\\
Visitation (2 Jul / 31 May) & Song 2:8--14; Luke 1:39--47; Secret: \textit{matris integritatem... sacravit} & Zeph 3:14--18 (or Rom 12:9--16b); Isa 12; Luke 1:39--56 & Consecrated virginal integrity (1962); Daughter of Zion bearing the Lord in her midst (postconciliar)\\
\end{fourstable}

\subsubsection{The Litany of Loreto and Ineffabilis Deus (1854)}

In the daily Marian devotion of the faithful, the Litany of Loreto---recorded in \textit{The Raccolta} (1910, No.\ 189, noting Pope Pius VII's indulgence decree of 30 September 1817)---places \textit{Foederis arca}, ``Ark of the covenant,'' at the architectural summit of a tenfold cluster of biblical and sanctuary titles:

\begin{itemize}
\item \textit{Vas spirituale} --- Spiritual Vessel;
\item \textit{Vas honorabile} --- Vessel of honour;
\item \textit{Vas insigne devotionis} --- Vessel of singular devotion;
\item \textit{Rosa mystica} --- Mystical Rose;
\item \textit{Turris Davidica} --- Tower of David;
\item \textit{Turris eburnea} --- Tower of ivory;
\item \textit{Domus aurea} --- House of gold;
\item \textit{Foederis arca} --- Ark of the covenant;
\item \textit{Janua caeli} --- Gate of heaven;
\item \textit{Stella matutina} --- Morning star.
\end{itemize}

The inner triad moves steadily inward and upward: \textit{Domus aurea} recalls the gold-overlaid interior of Solomon's Temple (1 Kgs 6:20--22); \textit{Foederis arca} passes behind the veil to the golden chest beneath the wings of the cherubim; and \textit{Janua caeli} opens that sanctuary onto heaven itself, where St.\ John beholds the temple of God opened and the Ark of his testament within (Rev 11:19).

When Pope Pius IX solemnly defined the dogma of the Immaculate Conception in the Apostolic Constitution \textit{Ineffabilis Deus} (8 December 1854), he gathered this entire patristic and liturgical stream into the teaching of the supreme magisterium. Recalling how the Fathers read Old Testament figures of Mary's spotless grace, Pius IX writes:

\begin{quote}
This illustrious and singular triumph of the Virgin, and her most exalted innocence, purity, and holiness, Her freedom from all stain of sin, and ineffable abundance of all heavenly graces, virtues, and privileges, the same fathers beheld in that ark of Noah which, divinely appointed, escaped safe and sound, from the common shipwreck of the whole world ... also in that most august temple of God, which, shining with divine splendor, is filled with the glory of God ... To describe the same totality, as it were, of Divine gifts, and the original integrity of the Virgin, of whom Jesus was born, the same Fathers, using the eloquence of the Prophets, celebrate the august Virgin as the spotless dove, the holy Jerusalem, the exalted throne of God, the ark and house of sanctification (\textit{arcam sanctificationis}), which Eternal Wisdom built for itself.
\end{quote}

A few paragraphs later, Pius IX reaches directly back to Bezalel, the Spirit-filled master craftsman who built the wilderness Tabernacle and the Ark of the Covenant in Exodus 31 and 37:

\begin{quote}
They have testified that the flesh of the Virgin, taken from Adam, did not admit the stains of Adam, and on this account that the Most Blessed Virgin was the tabernacle created by God himself formed by the Holy Spirit, truly enriched with purple which that new Beseleel made, adorned and woven with gold.
\end{quote}

In the original Latin (\textit{tabernaculum esse ab ipso Deo creatum, Spiritu Sancto formatum, et purpureae revera operae, quod novus ille Beseleel auro intextum variumque effinxit}), the Holy Spirit himself is the new Bezalel, crafting a living sanctuary from human flesh so that the Son of God might dwell among men.
